\documentclass[10pt,a4paper]{amsart}
\usepackage[margin=19mm]{geometry}
\usepackage{amsmath,amssymb,mathtools}
\usepackage[scr=rsfs,cal=euler]{mathalfa}
\usepackage{xfrac}
\usepackage[unicode,hidelinks,bookmarks=false]{hyperref}
\newcommand{\ie}{i.e.\ }
\newcommand{\cf}{cf.\ }
\newcommand{\resp}{respectively\ }
\newcommand{\Subsection}{\subsection}
\providecommand{\nobreakdash}{\nobreak\hbox{-}}
\setlength{\emergencystretch}{2em}
\multlinegap0pt
\usepackage{amssymb}
\usepackage{mathtools}
\newtheorem{assumption}{Assumption}
\newtheorem{lemma}{Lemma}
\def\a{\alpha}
\def\la{\langle}
\def\ra{\rangle}
\def\re{\mathbb{R}}
\title{Randomized Feasibility Methods for Constrained Optimization with Adaptive Step Sizes}
\author{Abhishek Chakraborty, Angelia Nedi\'c}
\date{Source: arXiv:2601.20076v2 (CC BY 4.0)}
\begin{document}
\maketitle

\noindent\emph{Source and licence.} Randomized Feasibility Methods for Constrained Optimization with Adaptive Step Sizes. Version arXiv:2601.20076v2 (CC BY 4.0), distributed by arXiv under the Creative Commons Attribution 4.0 International licence, http://creativecommons.org/licenses/by/4.0/, as recorded in the arXiv licence metadata for this exact version and shown on the abstract page https://arxiv.org/abs/2601.20076v2. Full licence notice: \texttt{licenses/arxiv-2601.20076-CC-BY-4.0.txt}.\par\smallskip

\noindent\textit{Editorial scope.} Counted unit: the article's lemma lem-elem-grad (source0.tex lines 1103--1110). The article's complete original proof of it is delivered with it (source0.tex lines 1112--1134, one \texttt{proof} environment of 23 lines). No mathematics is written, repaired or re-derived: the statement and the proof are the article's own lines, in the article's own notation, and no retained line is substituted or re-flowed.  Retained uncounted as the source-local context the proof needs: lines 385--390; lines 393--399.  Nothing was omitted from any retained span, so the package carries no omission record.  No retained line contains a citation, so the wrapper carries no bibliography and no bibliography entry.  No retained line contains a figure environment, an image-inclusion command or drawing code, so no image is delivered and the package declares no asset dependency.  Editorial formatting only, in addition: an AMS \texttt{amsart} preamble that restates the article's own declarations and macros used by the retained lines, namely the environments \texttt{assumption}, \texttt{lemma} on separate counters, together with the standard packages those lines need.  The retained environments print in this document as Assumption 1, Lemma 1 in order of appearance, whereas the article prints the same items with the numbering of its own full document.  The complete native source of the article as distributed in the arXiv e-print for this exact version is bundled unchanged as \texttt{sources/193503/source0.tex}.
\begin{assumption}\label{asum_lipschitz}
    The function $f:\re^n\to\re$ has Lipschitz continuous gradients over $Y$ with a constant $L>0$, i.e.,
   \begin{align}
       f(y)-f(x)\le \la \nabla f(x),y-x\ra+\frac{L}{2}\|y-x\|^2 \;\; \forall x,y\in Y. \nonumber
   \end{align}
\end{assumption}

\begin{assumption}\label{asum_strong_convex}
    The function $f:\re^n\to\re$ is strongly convex
   over $Y$ with the constant $\mu>0$, i.e.,
   \begin{align}
       f(y)-f(x)\ge \la \nabla f(x),y-x\ra+\frac{\mu}{2}\|y-x\|^2 \;\; \forall x, y \in Y . \nonumber
   \end{align}
\end{assumption}

\begin{lemma}\label{lem-elem-grad} 
     Let $Z$ be a closed convex set, and let 
     $f(\cdot)$ be a continuously differentiable function. Consider an update of the following form:
     \[w=\Pi_Z[x-\a\nabla f(x)],\]
     where $\a>0$ and $x\in\re^n$. If the function $f(\cdot)$ is strongly convex and has Lipschitz continuous gradients (Assumptions~\ref{asum_lipschitz} and~\ref{asum_strong_convex} hold), then 
     we have for all $z\in Z$,
     \[\|w - z\|^2 \leq (1-\a \mu) \| x -z \|^2 - (1 - \a L) \| w - x\|^2 - 2 \a (f(w) - f(z)).\]
     \end{lemma}

 \begin{proof}
   From the definition of $w$ and the non-expansiveness property of the projection,  we have for any $z \in Z$,
    \begin{align}\label{norm_diff11}
        \|w - z\|^2 &= \| \Pi_Z[x - \a \nabla f(x)] - z\|^2\cr
        &\leq \| x - \a \nabla f(x) - z\|^2 - \| w - (x - \a \nabla f(x)) \|^2 \cr
        & = \| x -z \|^2 - \| w - x \|^2 
        - 2 \a\la \nabla f(x), x -z \ra 
        - 2 \a \la \nabla f(x), w - x \ra.
    \end{align}
    We next bound the inner product terms on the right hand side of relation~\eqref{norm_diff11}. By the strong convexity of $f(\cdot)$ (Assumption~\ref{asum_strong_convex}),
    we have
    \begin{align}\label{third_term}
        - 2 \a \la \nabla f(x), x -z \ra \leq - 2 \a (f(x) - f(z)) - \a \mu \| x - z \|^2 .
    \end{align}
    Using the Lipschitz continuity of $\nabla f(\cdot)$ (Assumption~\ref{asum_lipschitz}), the last term on the right hand side of relation~\eqref{norm_diff11} can be bounded as follows:
    \begin{align}\label{fourth_term}
        - 2 \a \la \nabla f(x), w - x \ra 
        \leq 2 \a (f(x) - f(w)) + \a_k L \|w - x \|^2 .
    \end{align}
    Combining the estimates in relations~\eqref{third_term} and \eqref{fourth_term} with relation~\eqref{norm_diff11}, we obtain for any $z\in Z$,
    \[\|w - z\|^2 \leq (1-\a \mu) \| x -z\|^2 - (1 - \a L) \| w - x \|^2 - 2 \a (f(w) - f(z)),\]
    thus showing the stated relation.
\end{proof}

\end{document}
